\chapter{MTS et MTO : principes}\label{ch:mts-mto}

Une stratégie de production répond à une question : à quel moment la production est-elle
décidée, par rapport à la demande client ?

La stratégie MTS (\textit{Make-To-Stock}) produit par anticipation vers un stock de produits
finis. La demande est servie depuis ce stock, et la production ne dépend pas d'une commande
précise. La stratégie MTO (\textit{Make-To-Order}) produit à partir d'une commande client.
La nomenclature de la commande fixe les matières nécessaires, et la production ne démarre
qu'une fois ces matières disponibles pour cette commande. \cref{tab:mts-mto} résume les
différences, et \cref{fig:mts-mto} les chemins correspondants. Aucune de ces stratégies
n'est supérieure dans l'absolu : le choix dépend de la variabilité de la demande, du coût de
stockage et de la tolérance du client au délai.

\begin{table}[H]
  \centering
  \caption{Comparaison des stratégies MTS et MTO dans SCONTO-SVU.}\label{tab:mts-mto}
  \begin{tabularx}{\textwidth}{@{}>{\bfseries}L{3.2cm}YY@{}}
    \toprule
    Critère & MTS (\textit{Make-To-Stock}) & MTO (\textit{Make-To-Order}) \\
    \midrule
    Déclencheur de production & seuil de stock de produits finis (REAPPRO) & commande client \\
    Source de la demande servie & stock de produits finis & production dédiée à la commande \\
    Besoin matière & calculé sur la production anticipée & calculé par commande, selon la nomenclature \\
    Délai client & court si le stock couvre la commande & temps de production et d'approvisionnement \\
    Risque principal & immobilisation de stock & engagement de matière non disponible \\
    \bottomrule
  \end{tabularx}
\end{table}

\begin{figure}[H]
  \centering
  \begin{tikzpicture}[font=\small, >=Stealth, node distance=0pt,
    st/.style={draw=ink, rounded corners=3pt, minimum width=2.2cm, minimum height=0.9cm, align=center, fill=white}]
    % MTS : production anticipee vers le stock
    \node[font=\bfseries] at (0,3.4) {MTS : production anticipée};
    \node[st, fill=mtsc!18] (p1) at (0,2.4) {Production\\prévisionnelle};
    \node[st, fill=mtsc!18] (s1) at (3.2,2.4) {Stock de\\produits finis};
    \node[st] (c1) at (6.4,2.4) {Commande\\client};
    \node[st] (d1) at (9.6,2.4) {Livraison\\immédiate};
    \draw[->, thick] (p1) -- (s1);
    \draw[->, thick] (s1) -- (c1);
    \draw[->, thick] (c1) -- (d1);
    % MTO : production declenchee par la commande
    \node[font=\bfseries] at (0,0.9) {MTO : production à la commande};
    \node[st, fill=mtoc!22] (c2) at (0,-0.1) {Commande\\client};
    \node[st, fill=mtoc!22] (n2) at (3.2,-0.1) {Nomenclature\\et besoin matière};
    \node[st, fill=mtoc!22] (m2) at (6.4,-0.1) {Make\\(production)};
    \node[st] (d2) at (9.6,-0.1) {Deliver\\(livraison)};
    \draw[->, thick] (c2) -- (n2);
    \draw[->, thick] (n2) -- (m2);
    \draw[->, thick] (m2) -- (d2);
  \end{tikzpicture}
  \caption{Comparaison des chemins MTS (haut) et MTO (bas). En MTS, la commande est servie
    depuis un stock produit ; en MTO, elle déclenche sa propre production.}
  \label{fig:mts-mto}
\end{figure}


\section{Le chemin d'une commande MTO}

Pour une commande MTO, la chaîne logique est la suivante. Le client émet une commande. Le
système lit la nomenclature du produit, c'est-à-dire la liste des matières et des quantités
par unité produite, puis calcule le besoin matière de la commande. Si ce besoin n'est pas
couvert, il déclenche un approvisionnement. Quand les matières sont disponibles, la
production (Make) démarre. Les unités circulent sur les postes, consomment la matière à
l'étape concernée, puis sont livrées (Deliver). La commande est clôturée à la date de
livraison, avec le statut \Etat{SERVIE} si elle est livrée dans le délai promis, et
\Etat{EN\_RETARD} sinon.

Cette chaîne suppose implicitement que la disponibilité matière est une quantité libre,
qui n'est pas partagée avec d'autres commandes. Le chapitre suivant montre que cette
hypothèse ne tient plus dès que plusieurs commandes existent.

\section{Le problème de l'engagement}

Imaginons un stock de 200 unités de GPL en vrac. La commande \jv{CMD\_1} demande 125 unités.
Si le système lit le stock physique seul, la commande voit 200 unités disponibles. Mais si
\jv{CMD\_1} est acceptée et engage 125 unités, il ne reste que 75 unités réellement libres,
même si les 200 unités sont physiquement présentes. Une deuxième commande demandant 125
unités ne doit donc pas pouvoir consommer ce qui appartient déjà à \jv{CMD\_1}. Cette
distinction entre stock physique et stock libre est formalisée au chapitre suivant.
